
Reinforcement learning from execution feedback (RLEF) is known to improve code language models over supervised fine-tuning (SFT), but the mechanism and the importance of reward design remain incompletely characterized, particularly at hard benchmark difficulty levels. This paper presents a controlled, reproducible study comparing RLEF and SFT on DeepSeek-Coder-7B trained on the APPS dataset and evaluated across five benchmark difficulty levels --- from HumanEval to LiveCodeBench-Hard. The central finding is that RLEF's advantage over SFT is concentrated at the hardest evaluation level (LiveCodeBench-Hard, $\Delta =+0.18$, from smoke-scale proxy evaluation) and follows a statistically significant positive-ordered trend across difficulty levels (Jonckheere-Terpstra z=+56.10, $p\approx 0)$. Critically, this advantage holds regardless of reward formulation: fraction-of-tests and binary RLEF rewards produce statistically equivalent performance ($\Delta _Fraction - \Delta _Binary$ = +0.0072, p=0.552, 95\% CI [$-0.075$, +0.089]), consistent with recent literature reporting reward formulation convergence under GRPO. A secondary finding refines the standard mechanistic explanation: despite 85.32\% reference solution coverage in APPS competition problems, SFT achieves 0.0 pass@1 on LiveCodeBench-Hard, indicating a generalization void rather than a dataset void --- though the precise causal mechanism remains under investigation. All quantitative $\Delta$ values are derived from smoke-scale proxy evaluation (500 APPS samples, 62 GRPO steps, N=50 per benchmark); full-scale confirmation is deferred to future work. The training infrastructure and analysis code are fully open-source.

